\section{问题重述与分析}

\subsection{研究背景与任务要求}

赛题提供一份商用啤酒风味数据集，包含八份附件。
附件一记录 250 款啤酒的 231 个化学与理化特征，附件二给出这些特征两两之间的相关系数矩阵。
附件三定义感官记分卡，附件四给出 250 款啤酒在 50 个感官维度上的标准化品评小组均值。
附件八补充部分化合物的校准浓度、气味阈值、来源标签与缺陷描述。
这些附件对应 Schreurs 等人在 \emph{Nature Communications} 发表的数据记录，作者在同一数据记录的新版中另行公布了逐品评员、逐啤酒的原始 0 到 6 分评分，覆盖相同的 250 款啤酒，是三个问题的感官数据来源。

赛题要求完成三项任务。
问题一以化学数据预测啤酒的感官评分，问题二找出每个感官维度的关键化合物，问题三生成指定风格的一款啤酒在各化学特征上的取值向量，即化合物谱，并进行跨风格迁移。

\subsection{问题分析}

\subsubsection{评分预测}

问题一从化学特征预测每款啤酒在每个感官维度上的品评等级。
对啤酒 $b$（$b=1,\dots,250$）与感官维度 $d$（$d=1,\dots,50$），把该啤酒该维度的全部原始评分记录取算术平均得到 $\bar x_{bd}$，再截断到 0 到 6 并四舍五入，得到等级 $G_{bd}$。
交付物为每款啤酒每个维度的预测等级、等级概率与 90\% 等级区间。
难点在于样本量小、特征维数高：附件一的化学数据有缺失值，待评分的新啤酒没有任何品评记录，评分只能由化学数据预测；感官评分还含品评员之间的系统差异与随机噪声，同一天品评的啤酒共享出席名单与品评情境，不能当作相互独立的样本。
因此既要在训练折内处理化学缺失，又要用按日期组分组的嵌套交叉验证控制同日信息泄漏。

\subsubsection{化合物识别}

问题二的目标是找出每个感官维度上关联最强的三个化合物。
第 1 位取与维度秩相关最强的化合物，第 2、3 位在控制已选化合物之后取秩偏相关最强、且不属于已选化合物所在共表达模块的化合物。
这一目标把关键化合物定义为关联结果，而不是因果结果，附件中没有干预数据，化学合理性在结果产生后另行解释。
难点在于附件二只在全部 250 款啤酒上估计两两相关，风格之间的配方差异与风格内的共变混在一起；同一模块内的化合物高度相关，边际关联最强的几个化合物往往来自同一模块，直接取前 3 会重复选择同一条信息。
因此需要先用模拟还原实验确定选择规则，再用日期组重抽量化入选频率。

\subsubsection{风格生成}

问题三的估计目标有三项：风格条件分布、感官画像与迁移映射。
风格条件分布是给定细分风格后 225 个生成列的联合分布，感官画像是由该分布生成的化合物谱经问题一模型得到的 50 维等级分布，迁移映射是把一种风格的化合物谱改为另一种风格、同时满足化学层与感官层两层缺陷约束。
难点在于 22 个细分风格中每个只有 3 到 31 款啤酒，少于 225 个生成列，风格专属的完整协方差不可估计；生成结果既要落在观测范围内，又要让风格判别器无法区分生成样本与真实啤酒，还不能只是复制训练啤酒；迁移时题面要求缺陷化合物不恶化，需要写成可检验的约束，并在潜空间与原始尺度两处复查。

三项任务共享同一批啤酒与化学特征，按问题一、问题二、问题三的顺序衔接，前两问的产物作为第三问的输入。
问题一建立化学谱到 50 维感官评分分布的映射；问题二借助共表达模块识别每个感官维度的关键化合物候选；问题三建立以风格为条件的生成模型，调用问题一的评分映射形成感官画像，并在化学与感官两层约束下完成跨风格迁移。
